\section{Experiments}\label{sec:exp}
\paragraph{Protocol.} Every number comes from Monte-Carlo replicates with a known truth. \emph{Real-data cells:} the seven \texttt{ppi-python} example sets that contain $(Y,\hat Y)$ for all records (Galaxies $P{=}16{,}743$, Forest $1{,}596$, AlphaFold $10{,}802$, Gene expression $61{,}150$, Census income $380{,}091$, Census health insurance $318{,}215$, Census education $318{,}215$ records; $P$ is the population size) are each treated as the population; a replicate draws $n$ labeled and $N$ unlabeled records i.i.d.\ \emph{with replacement} and the truth is the population mean of $Y$. The predictor is the one shipped with the data, never refit. (Records repeat when $n+N$ is not small relative to $P$, most visibly for Forest; the DGP is still i.i.d.\ from the empirical distribution, but the efficiency bound of Prop.~\ref{prop:bound} for that DGP is then not the one we compute from cross-validated fits, so we flag Forest at $n\ge1000$.) \emph{Synthetic cells:} (a) $f\sim N(0,1)$, $Y=f+c\,(f^2-1)/\sqrt2+\varepsilon$, $\varepsilon\sim N(0,1)$, so $s=c^2/(1+c^2)$ is known exactly; (b) a classifier that outputs logits $f\sim N(\mu,\tau^2)$ with $Y\sim\mathrm{Bernoulli}(\mathrm{expit}(f))$ (so $\mathbb E[Y|f]=\mathrm{expit}(f)$ is non-affine), with $(\mu,\tau)=(-1,2)$ ($s=0.06$) and $(-1,4)$ ($s=0.23$). \emph{Methods:} classical; PPI ($\lambda{=}1$) and PPI++ from \texttt{ppi-python} 0.2.3 (which clips $\hat\lambda$ to $[0,1]$); PPI++ with the unclipped plug-in $\hat\lambda=\widehat{\mathrm{Cov}}(Y,f)/((1+n/N)\widehat{\mathrm{Var}}(f))$, which is the best member of the PPI family and is our reference; and cross-fitted lin / bin10 / iso / spl / auto (Sec.~\ref{sec:method}, $K=5$). All intervals are 95\% Wald-type. \emph{Metrics:} coverage with Wilson intervals; width reported as the paired per-replicate ratio to the reference (mean $\pm1.96$ s.e.\ over replicates, so values that differ from 1 by more than the $\pm$ are not noise); the closed-form ratio $\sqrt{V_{\rm eff}/V_{\rm PPI++}}$ from Props.~\ref{prop:ppi}--\ref{prop:bound} (``bound'') computed from the population (real data: 5-fold cross-validated isotonic/spline fit of $\mathbb E[Y|f]$, so it is itself an estimate, and is capped by the linear fit). Replicates: 2000 per synthetic cell, 1000 per real-data cell. Seeds are derived from (experiment cell, replicate index), so the same labeled/unlabeled draws are shared across methods within a cell.

\subsection{Synthetic data: the bound is tight, and recalibration reaches it}
Figure~\ref{fig:theory} varies the nonlinear share $s$ at $n=500$. The cross-fitted linear interval tracks PPI++ (within 1.2\%, the small price of cross-fitting), while spline/auto intervals follow the efficiency bound: for $N{=}50{,}000$ the width ratio to PPI++ is 0.47 at $s=0.8$, 0.72 at $s=0.5$, 0.91 at $s=0.2$ against the bound 0.46, 0.71, 0.90. Over the 20 cells the empirical ratio of spline/auto exceeds the bound by $0.010\pm0.003$ (mean $\pm$ s.d. over cells; maximum 0.015), which is the cost of estimating the calibration curve. Gains vanish when $N$ is small ($N{=}100$: at most 5.5\% at $s=0.8$), as Cor.~\ref{cor:gap} predicts. Isotonic regression falls short here because $c(f^2-1)$ is non-monotone: its limit $g^*\neq g$, so the interval stays valid but longer (0.67 vs 0.47), as Prop.~\ref{prop:valid} allows. Coverage of auto was between 94.2\% and 95.9\% in all 20 cells (the two cells where several methods sit at 93.4--94.0\% are shared by PPI and PPI++ because the replicates are shared; Table~\ref{tab:valid}).

\begin{figure}[t]\centering
\includegraphics[width=\linewidth]{../figs/fig_theory.pdf}
\caption{Width of cross-fitted intervals relative to PPI++ (unclipped $\lambda$) versus nonlinear share $s$ of the calibration curve ($n=500$; synthetic quadratic DGP; 2000 replicates; bars: $\pm1.96$ s.e.). Dashed: closed-form bound $\sqrt{V_{\rm eff}/V_{\rm PPI++}}$.}\label{fig:theory}
\end{figure}

\paragraph{Small samples (Fig.~\ref{fig:smalln}).} The price of flexibility is paid at small $n$. With a linear truth ($c=0$), auto is 14.5\%, 7.1\%, 3.4\%, 0.6\% wider than PPI++ at $n=25,50,100,500$; spline alone is 65\% wider at $n=25$. With a quadratic truth ($c=1$) auto breaks even at about $n=30$ and gains $10.1\%$ ($n=50$), 20\% ($n=100$), 27\% ($n=500$). Wald coverage at $n=25$ is 91--95\% across methods (classical 92.7--93.0\%); the cross-fitted intervals are not worse and, at $n=25$--$50$, better than ppi\_py's PPI++ (e.g.\ $c=1$, $n=50$: 93.2\% auto vs 91.8\% PPI++, 92.0\% for the $t$-version of unclipped PPI++ vs 93.6\% for auto). A Student-$t$ calibration of the half-widths with $n-1$ d.f.\ narrows but does not close this gap for PPI++ (92.0\% at $n=50$).

\begin{figure}[t]\centering
\includegraphics[width=.82\linewidth]{../figs/fig_smalln.pdf}
\caption{Small-$n$ behaviour ($N=20{,}000$, quadratic synthetic DGP, 2000 replicates). Top: width relative to PPI++ (unclipped). Bottom: coverage with Wilson 95\% intervals.}\label{fig:smalln}
\end{figure}

\subsection{Real data with the shipped predictors: essentially nothing to gain}\label{sec:exp-real}
Table~\ref{tab:real} (and Appendix Figs./Tables) shows the central negative result. The nonlinear share of the calibration curve is at most $2.5\%$ (Gene expression) and $\le0.5\%$ for the other six data sets; the implied bound on the achievable width reduction over PPI++ at $n=500$, $N=20{,}000$ is 0.9\% (Gene expression), 0.2\% (Census income) and $\le0.1\%$ elsewhere. The cross-fitted auto interval does not realise even that: at $n=500$ it is 0.5--2.1\% \emph{wider} than PPI++ on six data sets, and 0.6\% narrower on Gene expression; at $n=2000$ all six are within 0.6\% of PPI++ (the 1\% ``gains'' on Forest at $n\ge1000$ are the resampling artifact flagged above); at $n=50$ auto is 3--19\% wider except on Gene expression ($-2.4\%$). At $n\ge500$ the coverage of every method lies between 91.5\% and 96.8\% (mean 94.5--95.2\%); the lowest cells (92.2--92.8\%, Census income with $N=5000$) are shared by all methods including the classical interval. At smaller $n$ the Census-income coverage is 82--92\% for $n\le200$ for every method (heavy-tailed income), so these shortfalls come from the Wald approximation and not from recalibration.

Two other facts matter for practice. (i) The headroom \emph{relative to the classical interval} is large (width 1.2--3.2 times higher), so using a predictor is worthwhile; the headroom beyond PPI++ is not. (ii) The choice \emph{within} the PPI family matters more than recalibration: PPI with $\lambda=1$ is 4.9\% wider than PPI++ on Gene expression ($\hat\beta=1.15$) and 1.8\% on Census income ($\hat\beta=1.38$), and ppi\_py's PPI++ shows the same loss because its $\hat\lambda$ is clipped at 1 while the unclipped optimum is above 1 (bound ratio of the clipped version to the unclipped one: 1.056 and 1.018).

\begin{table}[t]\centering\small
\caption{Real data, shipped predictors, $n=500$, $N=20{,}000$, 1000 replicates. corr: correlation of $Y$ and $\hat Y$; $\hat\beta$ slope of $Y$ on $\hat Y$; $s$: cross-validated nonlinear share of residual variance; widths are ratios to PPI++ with unclipped $\hat\lambda$ ($\pm$: 1.96 s.e.); ``bound'' is $\sqrt{V_{\rm eff}/V_{\rm PPI++}}$ estimated from the population. Forest has only 1596 records ($n/P=0.31$).}\label{tab:real}
\resizebox{\linewidth}{!}{\begin{tabular}{lrrr rrrr rrrr}
\toprule
 & & & & \multicolumn{4}{c}{coverage (\%)} & \multicolumn{4}{c}{width / PPI++ (unclipped)}\\
\cmidrule(lr){5-8}\cmidrule(lr){9-12}
data set & corr & $\hat\beta$ & $s$ & classical & PPI & PPI++ & auto & classical & PPI & auto & bound\\
\midrule
Galaxies & 0.74 & 0.98 & 0.0\% & 95.3 & 95.1 & 94.7 & 95.2 & 1.47 & 1.00 & 1.007$\;{\scriptstyle\pm0.000}$ & 1.000\\
Forest & 0.54 & 0.90 & 0.0\% & 96.1 & 95.3 & 95.1 & 95.1 & 1.18 & 1.00 & 1.005$\;{\scriptstyle\pm0.001}$ & 1.000\\
AlphaFold & 0.82 & 1.03 & 0.2\% & 95.0 & 94.8 & 94.3 & 94.2 & 1.69 & 0.99 & 1.005$\;{\scriptstyle\pm0.001}$ & 0.999\\
Gene expression & 0.96 & 1.15 & 2.5\% & 94.1 & 93.9 & 94.0 & 93.8 & 3.24 & 1.05 & 0.994$\;{\scriptstyle\pm0.001}$ & 0.991\\
Census income & 0.61 & 1.38 & 0.5\% & 94.2 & 93.5 & 93.1 & 93.4 & 1.26 & 1.02 & 1.021$\;{\scriptstyle\pm0.002}$ & 0.998\\
Census health ins. & 0.71 & 0.99 & 0.1\% & 94.2 & 94.7 & 94.8 & 95.0 & 1.40 & 1.00 & 1.008$\;{\scriptstyle\pm0.000}$ & 1.000\\
Census education & 0.63 & 0.98 & 0.0\% & 93.9 & 94.7 & 94.1 & 94.5 & 1.28 & 1.00 & 1.011$\;{\scriptstyle\pm0.001}$ & 1.000\\
\bottomrule
\end{tabular}}
\end{table}

\subsection{Distorted predictor outputs: where recalibration pays}\label{sec:exp-distort}
Classifiers often emit scores that are monotone but not affine in the probability of the label (logits, over-confident or squashed scores \citealp{niculescu2005predicting,guo2017calibration}). To measure what that costs, we apply a fixed monotone map to the shipped predictor outputs (logit of the probability score; cubed score; $\exp(\hat y/4)$ for Gene expression; $\sqrt{\cdot}$ for Census income) and keep everything else. These are \emph{constructed} cases, chosen to represent the miscalibration mechanism, not evidence of how common it is. Table~\ref{tab:distort} shows that the width reduction of auto/iso relative to PPI++ ranges from $-2\%$ to $-27\%$, matching the bound within about 2 percentage points ($n=1000$), with nominal coverage; plain PPI ($\lambda{=}1$) becomes 4.9--9.5 times wider than PPI++ for logit scores (and 4.6 times for the exponential distortion) because the scores are no longer on the scale of $Y$. With logit-scale synthetic scores of moderate sharpness ($s=0.06$) the gain is only about 2\%, with sharper logits ($s=0.23$) it is 10\% at $n=500$ (bound 11.7\%) and 3\%, 8\%, 10\% at $n=50,100,200$ (Appendix Table~\ref{tab:logit}).

\begin{table}[t]\centering\small
\caption{Monotonically distorted predictor outputs, $n=1000$, $N=20{,}000$, 1000 replicates. Width ratios to PPI++ with unclipped $\hat\lambda$. The last column is the ratio of ppi\_py's clipped PPI++ to the unclipped one.}\label{tab:distort}
\resizebox{\linewidth}{!}{\begin{tabular}{l r rrrr rrr r}
\toprule
 & & \multicolumn{4}{c}{width / PPI++ (unclipped)} & \multicolumn{3}{c}{coverage (\%)} & \\
\cmidrule(lr){3-6}\cmidrule(lr){7-9}
predictor output & $s$ & PPI & iso & auto & bound & PPI++ & iso & auto & $\lambda^{\rm clip}$-PPI++\\
\midrule
Galaxies, logit & 9\% & 8.46 & 0.964$\;{\scriptstyle\pm0.001}$ & 0.963$\;{\scriptstyle\pm0.001}$ & 0.957 & 95.1 & 94.3 & 94.2 & 1.00\\
Galaxies, pow3 & 13\% & 1.00 & 0.945$\;{\scriptstyle\pm0.001}$ & 0.942$\;{\scriptstyle\pm0.001}$ & 0.937 & 95.1 & 95.2 & 96.0 & 1.00\\
AlphaFold, logit & 23\% & 9.50 & 0.900$\;{\scriptstyle\pm0.001}$ & 0.899$\;{\scriptstyle\pm0.001}$ & 0.892 & 94.7 & 95.0 & 95.0 & 1.00\\
Forest, logit & 5\% & 4.94 & 0.975$\;{\scriptstyle\pm0.001}$ & 0.977$\;{\scriptstyle\pm0.001}$ & 0.979 & 94.7 & 95.1 & 95.2 & 1.00\\
Gene expression, exp & 60\% & 4.62 & 0.738$\;{\scriptstyle\pm0.001}$ & 0.731$\;{\scriptstyle\pm0.001}$ & 0.731 & 94.8 & 95.4 & 95.3 & 1.00\\
Census health ins., logit & 8\% & 6.50 & 0.971$\;{\scriptstyle\pm0.001}$ & 0.967$\;{\scriptstyle\pm0.001}$ & 0.962 & 94.4 & 94.3 & 94.5 & 1.00\\
Census income, sqrt & 8\% & 1.20 & 0.987$\;{\scriptstyle\pm0.002}$ & 0.976$\;{\scriptstyle\pm0.001}$ & 0.961 & 94.6 & 95.5 & 94.9 & 1.20\\
\bottomrule
\end{tabular}}
\end{table}

\subsection{Binary predictions, distribution shift, and a fixed-sets application}\label{sec:exp-shift}
\paragraph{Binary predictor.} The plankton predictions are 0/1. Then $g$ is affine and Cor.~\ref{cor:gap}(ii) applies: auto/iso/lin are within 0.1\% of PPI++ in width for $n\ge500$ and 1.2\% narrower at $n=100$ (coverage in Table~\ref{tab:shift}, top). The same holds for the ballots application below.
\paragraph{Distribution shift.} The labeled plankton records are from 2013 and the unlabeled ones from 2014; the target is the 2014 mean. PPI's rectifier $Y-\hat Y$ is approximately stable across years, so PPI ($\lambda{=}1$) keeps 94\% coverage until $n=2000$ and degrades at $n=10{,}000$ (91\%), whereas every interval that estimates the calibration curve on 2013 labels (PPI++ with $\hat\lambda<1$, lin, iso, auto) collapses: auto covers 84\% at $n=500$, 9\% at $n=2000$, 0\% at $n=10{,}000$. Efficiency comes from leaning harder on the stationarity of $\mathbb E[Y|f]$, which is exactly what fails under shift. Estimating the calibration curve thus makes the interval \emph{less} robust than the PPI default; we have no remedy and make no claim about shift-robust variants.
\paragraph{Fixed-sets application.} For the ballots data ($n=1027$, $N=78{,}150$, binary predictions, no ground truth for the unlabeled records) the classical interval has width 0.0591, PPI 0.0133, PPI++ 0.0131 and all cross-fitted intervals 0.0132 (Appendix Table~\ref{tab:ballots}); with a binary predictor nothing below PPI++ is possible, in line with the theory. Here $\hat Y=1\Rightarrow Y=1$ for all labeled records, so the predictor is nearly perfect on one side.

\begin{table}[t]\centering\small
\caption{Coverage (\%) on plankton (1000 replicates, 95\% nominal). Top: labeled and unlabeled data from the same year (2013); bottom: labeled 2013, unlabeled 2014, target is the 2014 mean. $N=20{,}000$.}\label{tab:shift}
\resizebox{\linewidth}{!}{\begin{tabular}{ll rrrrr}
\toprule
setting & $n$ & classical & PPI & PPI++ & lin & auto\\
\midrule
no shift (2013 only) & 100 & 92.3 & 95.2 & 91.7 & 95.6 & 95.2\\
 & 500 & 94.3 & 95.7 & 94.2 & 94.8 & 94.7\\
 & 2000 & 94.8 & 94.5 & 94.0 & 94.3 & 94.5\\
 & 10000 & 94.0 & 95.1 & 94.6 & 94.7 & 94.7\\
\midrule
shift 2013$\to$2014 & 100 & 53.7 & 94.2 & 94.1 & 96.5 & 96.3\\
 & 500 & 0.4 & 94.8 & 94.8 & 84.2 & 84.4\\
 & 2000 & 0.0 & 94.2 & 94.2 & 9.1 & 8.9\\
 & 10000 & 0.0 & 91.2 & 0.5 & 0.0 & 0.0\\
\bottomrule
\end{tabular}}
\end{table}

\subsection{Side information beyond the predictor}\label{sec:exp-x}
For the three census data sets the features $X$ are available. Using $X$ in the calibration (effectively a new predictor trained on the labels) is outside the fixed-predictor constraint, but the bound shows how much it could give. Cross-validated gradient-boosting fits of $Y$ on $(f,X)$ lower the residual variance by at most 0.9\% (Census income; 0.1\% and 0.3\% for the other two) relative to a fit on $f$ alone (Appendix Table~\ref{tab:xgain}), since $f$ already summarizes $X$ (age/sex, income, etc.); there is no headroom here either. We did not run interval-level experiments for $X$-augmented calibration; the statement is about the population risk only.

\subsection{Validity summary}
Across all real and synthetic cells with $n\ge100$ (plankton no-shift cells included, shift cells excluded) Table~\ref{tab:valid} shows that no method is systematically invalid; coverage is significantly below 95\% (Wilson upper limit $<95\%$) in 12 of 157 cells for classical and the cross-fitted lin, 14 for iso/bin10/auto, but 22--26 for PPI, PPI++ (both versions); the worst cells ($\approx$ 87--91\%) are Census income at $n=100$ (all methods). Because the replicates are shared across methods these counts are strongly correlated, and part of the flags are chance (about 4 of 157 expected under exact 95\% coverage), so they should be read as ``the cross-fitted intervals are at least as well calibrated as the PPI family'', not as exact validity. Cells above 95\% are rare (1--7).

\begin{table}[t]\centering\small
\caption{Coverage over all cells with $n\ge100$ (real, distorted, synthetic and plankton no-shift; shift cells excluded).}\label{tab:valid}
\resizebox{\linewidth}{!}{\begin{tabular}{lrrrrr}
\toprule
method & cells & mean cov. (\%) & min cov. (\%) & \#cells sig.\ below 95 & \#cells sig.\ above 95\\
\midrule
classical & 157 & 94.64 & 90.9 & 12 & 1\\
PPI & 157 & 94.56 & 88.5 & 22 & 0\\
PPI++ (ppi\_py) & 157 & 94.51 & 87.7 & 24 & 2\\
PPI++ (unclipped) & 157 & 94.46 & 87.1 & 26 & 2\\
cross-fit linear & 157 & 94.80 & 88.9 & 12 & 2\\
cross-fit isotonic & 157 & 94.78 & 89.0 & 14 & 7\\
cross-fit 10 bins & 157 & 94.81 & 88.6 & 14 & 3\\
cross-fit spline & 157 & 94.77 & 88.4 & 13 & 2\\
cross-fit auto & 157 & 94.80 & 89.1 & 14 & 3\\
\bottomrule
\end{tabular}}
\end{table}
